% !TEX program = lualatex
\documentclass[aspectratio=43,professionalfonts,handout]{beamer}
\usepackage{beamer_template}

\newcommand{\LectureNum}{24}
\newcommand{\LectureTitle}{2倍角・半角の公式}
\newcommand{\TextPages}{p.\,160-162}
\newcommand{\NextTopic}{積を和・差に直す公式}
\newcommand{\NextPages}{p.\,163-164}
\newcommand{\CourseName}{基礎数学B}
\renewcommand{\TermName}{後期}

\begin{document}

% -----------------------------------------------
% スライド1: 表紙＋目標＋用語・定義
% -----------------------------------------------
\begin{frame}{\CourseName\quad 第\LectureNum 回「\LectureTitle 」}
  {\small \TermName\quad \TeacherName\quad ／\quad 教科書 \TextPages\quad
  \textbf{目標}：2倍角・半角の公式を導き、使いこなせる}

  \begin{block}{用語・定義}
  \begin{description}[labelwidth=5.5em]
    \item[\kw{$\sin 2\alpha = 2\sin\alpha\cos\alpha$}] $\sin 2\alpha = 2\sin\alpha\cos\alpha$
    \item[\kw{$\cos 2\alpha = \cos^2\alpha - \sin^2\alpha = 1 - 2\sin^2\alpha = 2\cos^2\alpha - 1$}]
      $\cos 2\alpha = \cos^2\alpha - \sin^2\alpha = 1 - 2\sin^2\alpha = 2\cos^2\alpha - 1$
    \item[\kw{$\tan 2\alpha = 2\tan\alpha/(1 - \tan^2\alpha)$}] $\tan 2\alpha = \dfrac{2\tan\alpha}{1-\tan^2\alpha}$
  \end{description}
  \end{block}
\end{frame}

% -----------------------------------------------
% スライド2: 公式・性質
% -----------------------------------------------
\begin{frame}{公式・性質}
  \begin{block}{}
  \begin{itemize}
    \item 半角：$\sin^2(\alpha/2) = (1 - \cos\alpha)/2$、$\cos^2(\alpha/2) = (1 + \cos\alpha)/2$
    \item 2倍角を使った方程式・不等式の解き方
  \end{itemize}
  \end{block}
\end{frame}

% -----------------------------------------------
% スライド3: 例1→問1
% -----------------------------------------------
\begin{frame}{例1→問1}
  \begin{exampleblock}{例2) p.177（2倍角で値を求める）}
    （板書で解説）
  \end{exampleblock}

  \vfill

  \begin{block}{問3) p.177\quad 自分で解いてみよう}
    （ノートに解くこと）
  \end{block}
\end{frame}

% -----------------------------------------------
% スライド4: 例2→問2
% -----------------------------------------------
\begin{frame}{例2→問2}
  \begin{exampleblock}{例3) p.178（三角関数の合成）}
    （板書で解説）
  \end{exampleblock}

  \vfill

  \begin{block}{問4) p.178\quad 自分で解いてみよう}
    （ノートに解くこと）
  \end{block}
\end{frame}

% -----------------------------------------------
% まとめ
% -----------------------------------------------
\begin{frame}{まとめ}
  \begin{itemize}
    \item 2倍角公式は加法定理で $\beta = \alpha$ とすれば導出できる
    \item $\cos 2\alpha = 1 - 2\sin^2\alpha$ から $\sin^2\alpha = \dfrac{1-\cos 2\alpha}{2}$（半角の公式）
    \item $\cos^2\alpha = \dfrac{1+\cos 2\alpha}{2}$ も半角の公式として使う
  \end{itemize}

  \vfill
  {\small 基本問題：321, 322, 323}
\end{frame}

\end{document}
